%% source: 2022-fa-redemption_midterm_01
%% tags: [binary search trees]
\begin{prob}
    Define the \textbf{largest gap} in a collection of numbers (also known as its \textbf{range}) to be
    greatest difference between two elements in the collection (in absolute value). For example, the
    largest gap in $\{4, 9, 5, 1, 6\}$ is 8 (between 1 and 9).

    Suppose a collection of $n$ numbers is stored in a \textbf{balanced} binary search tree. What is the
    time complexity required of an efficient algorithm to calculate the largest gap in the collection?
    State your answer as a function of $n$ in asymptotic notation.

    \begin{soln}
        $\Theta(\log n)$. To find the largest gap, we simply need to find the minimum and the maximum in 
        the balanced BST, which both take $\Theta(\log{n})$ time. Therefore adding them up, the time complexity for 
        calculating the largest gap is $\Theta(\log{n})$.
    \end{soln}

\end{prob}
